\subsubsection{Representación pentádica de la microfibra y de su incidencia}
\label{pi-estrella:desarrollo}

La disposición de las cinco regiones permite visualizar simultáneamente
la unidad del cociente ternario y la diferenciación que conserva el
catálogo. La aplicación relevante es la restricción
\[
 \Psi\big|_{\mathscr R_\pi^{\rm micro}}:
 \mathscr R_\pi^{\rm micro}\longrightarrow\{010211\},
 \qquad \Psi(U)=(-U)\bmod3.
\]
Sus cinco antecedentes son los vectores \(U_6\) de la tabla anterior.
La figura \ref{pi-estrella:figura} recupera su organización radial,
conservando el orden de coordenatización, las firmas y las multiplicidades del catálogo
construido en las secciones precedentes.

% Orden de carta y datos recuperados de fig_cinco_regiones_pi.tex.
\begin{figure}[H]
\centering
\begin{tikzpicture}[x=1cm,y=1cm,font=\small,>=Latex]
 \coordinate (a) at (90:2.2);
 \coordinate (b) at (18:2.2);
 \coordinate (c) at (-54:2.2);
 \coordinate (d) at (-126:2.2);
 \coordinate (e) at (162:2.2);
 \draw[black!35,line width=.45pt]
    (a)--(c)--(e)--(b)--(d)--cycle;
 \node[draw,fill=white,circle,minimum size=1.6cm,align=center,
    inner sep=3pt] (z) at (0,0) {$w_\pi$\\$010211$};
 \foreach \v in {a,b,c,d,e}{
    \fill (\v) circle (2pt);
    \draw[->,line width=.6pt] (\v)--(z);
 }
 \node[above=3mm,align=center] at (a)
   {$U_\pi^\perp$\\$555555$\\multiplicidad $432$};
 \node[right=3mm,align=left] at (b)
   {$U_{\pi,1}^{++}$\\$339555$\\multiplicidad $144$};
 \node[below right=2mm and 1mm,align=left] at (c)
   {$U_{\pi,2}^{++}$\\$393555$\\multiplicidad $144$};
 \node[below left=2mm and 1mm,align=right] at (d)
   {$U_{\pi,3}^{++}$\\$933555$\\multiplicidad $144$};
 \node[left=3mm,align=right] at (e)
   {$U_\pi^{--}$\\$933717$\\multiplicidad $144$};
 \node[align=center,font=\footnotesize] at (0,-3.55)
   {$1_\perp+3_{++}+1_{--}$\qquad $432+4\cdot144=1008$};
\end{tikzpicture}
\caption{Las cinco regiones de \(\pi\) en disposición pentádica.
Cada posición conserva su firma multiplicativa y su multiplicidad;
las flechas radiales representan la proyección común
\(\Psi(U)=010211\). El contorno estrellado organiza las cinco
posiciones de la coordenatización. Su realización en las cinco octadas de Witt
se efectúa mediante el alineamiento regional y la descomposición
\(4+1\) del texto.}
\label{pi-estrella:figura}
\end{figure}


La proyección al centro identifica las cinco imágenes ternarias; el
registro de las posiciones exteriores retiene la información regional
que hace posible la realización incidencial. La distribución
\(1_\perp+3_{++}+1_{--}\) especifica una región transversal, tres
regiones positivas y una negativa. Sus pesos normalizados son
\(3/7\) y \(1/7\) para cada una de las cuatro restantes: la
multiplicidad transversal es triple, aun cuando las cinco imágenes bajo
\(\Psi\) coincidan.

\Needspace{11\baselineskip}
El enlace con la estrella de Witt se formula mediante la bandera
\(B_K\subset H^-_\alpha\), el alineamiento \(\ell\) y la
elevación \(\iota_D\) ya construidos. Las cinco posiciones regionales
se realizan en las cinco octadas sobre \(\ell(B_K)\):
\[
\begin{aligned}
 U_\pi^\perp&\longmapsto O^\perp_{\ell(B_K)},\\
 U_\pi^{--}&\longmapsto\iota_D(H^-_\alpha),\\
 \{U_{\pi,1}^{++},U_{\pi,2}^{++},U_{\pi,3}^{++}\}
 &\longmapsto
 \{\iota_D(B_K\cup P):P\in\mathcal P_+\},\\
 \mathcal P_+&=\bigl\{\{1,3\},\{2,11\},\{8,12\}\bigr\}.
\end{aligned}
\]
La asignación de las tres posiciones positivas a los tres pares se
efectúa en una coordenatización ordenada. La fila transversal y la fila negativa
quedan distinguidas por sus predicados regionales y por la bandera.

Así se articulan las dos figuras: la figura
\ref{exc:fig-estrella} muestra las cuatro hexadas del pequeño diseño
de Witt sobre \(B_K\); la figura \ref{pi-estrella:figura} muestra
las cinco regiones que, bajo el alineamiento, ocupan las cuatro octadas
residuales y la octada transversal del diseño binario. La ley
\(4+1\), demostrada en la proposición
\ref{exc:residuo-binario}, especifica el cambio de incidencia. En esta
composición, \(K\) determina la cuaterna marcada y la firma de
\(\alpha\) distingue una de sus hexadas. Las regiones conservan,
además de su valor ternario común, la estructura necesaria para ese
transporte hacia la realización excepcional.
